% Options for packages loaded elsewhere
\PassOptionsToPackage{unicode}{hyperref}
\PassOptionsToPackage{hyphens}{url}
\PassOptionsToPackage{dvipsnames,svgnames,x11names}{xcolor}
\documentclass[
  11pt,
  letterpaper,
]{article}
\usepackage{xcolor}
\usepackage[margin=0.5in]{geometry}
\usepackage{amsmath,amssymb}
\setcounter{secnumdepth}{-\maxdimen} % remove section numbering
\usepackage{iftex}
\ifPDFTeX
  \usepackage[T1]{fontenc}
  \usepackage[utf8]{inputenc}
  \usepackage{textcomp} % provide euro and other symbols
\else % if luatex or xetex
  \usepackage{unicode-math} % this also loads fontspec
  \defaultfontfeatures{Scale=MatchLowercase}
  \defaultfontfeatures[\rmfamily]{Ligatures=TeX,Scale=1}
\fi
\usepackage{lmodern}
\ifPDFTeX\else
  % xetex/luatex font selection
  \setmainfont[]{Arial}
\fi
% Use upquote if available, for straight quotes in verbatim environments
\IfFileExists{upquote.sty}{\usepackage{upquote}}{}
\IfFileExists{microtype.sty}{% use microtype if available
  \usepackage[]{microtype}
  \UseMicrotypeSet[protrusion]{basicmath} % disable protrusion for tt fonts
}{}
\makeatletter
\@ifundefined{KOMAClassName}{% if non-KOMA class
  \IfFileExists{parskip.sty}{%
    \usepackage{parskip}
  }{% else
    \setlength{\parindent}{0pt}
    \setlength{\parskip}{6pt plus 2pt minus 1pt}}
}{% if KOMA class
  \KOMAoptions{parskip=half}}
\makeatother
\usepackage{longtable,booktabs,array}
\newcounter{none} % for unnumbered tables
\usepackage{calc} % for calculating minipage widths
% Correct order of tables after \paragraph or \subparagraph
\usepackage{etoolbox}
\makeatletter
\patchcmd\longtable{\par}{\if@noskipsec\mbox{}\fi\par}{}{}
\makeatother
% Allow footnotes in longtable head/foot
\IfFileExists{footnotehyper.sty}{\usepackage{footnotehyper}}{\usepackage{footnote}}
\makesavenoteenv{longtable}
\setlength{\emergencystretch}{3em} % prevent overfull lines
\providecommand{\tightlist}{%
  \setlength{\itemsep}{0pt}\setlength{\parskip}{0pt}}
\usepackage{bookmark}
\IfFileExists{xurl.sty}{\usepackage{xurl}}{} % add URL line breaks if available
\urlstyle{same}
\hypersetup{
  colorlinks=true,
  linkcolor={black},
  filecolor={black},
  citecolor={black},
  urlcolor={black},
  pdfcreator={LaTeX}}

\author{}
\date{}

\begin{document}

\section{Response to Reviewers}\label{response-to-reviewers}

\textbf{Manuscript:} Validated Synthetic Patient Generation for Small Longitudinal Cohorts: Coagulation Dynamics Across Pregnancy\\
\textbf{Journal:} npj Systems Biology and Applications\\
\textbf{Decision:} Major revision\\
\textbf{Date:} 7 August 2026

\begin{center}\rule{0.5\linewidth}{0.5pt}\end{center}

\subsection{Note on this document}\label{note-on-this-document}

Each reviewer comment is reproduced below. The comments are paraphrased, but all key points are retained. Each comment is followed by our response and the location of the corresponding manuscript change. Added or revised manuscript text is color-coded by reviewer:

\begin{itemize}
\tightlist
\item
  \textbf{Reviewer 1 edits: blue} (\texttt{\textbackslash{}rone\{...\}})
\item
  \textbf{Reviewer 2 edits: red} (\texttt{\textbackslash{}rtwo\{...\}})
\item
  \textbf{Edits addressing both reviewers: violet} (\texttt{\textbackslash{}rboth\{...\}})
\end{itemize}

A \texttt{\textbackslash{}reviewmode} setting in the preamble changes all colored text to black for the final version.

Reviewer 1\textquotesingle s original report numbered the comments 1 through 9 and then 11 through 13, with no comment 10. We kept every comment in its original order and renumbered the final three as R1.10, R1.11, and R1.12. These correspond to comments 11, 12, and 13 in the original report. No comment was omitted.

\begin{center}\rule{0.5\linewidth}{0.5pt}\end{center}

\subsection{General response to reviewers}\label{general-response-to-reviewers}

We thank both reviewers for their careful and constructive comments. Both recognized the value of combining statistical and mechanistic validation for very small longitudinal cohorts. They also identified four main concerns: generalizability beyond the single cohort of 23 patients, the separate contributions of principal-component analysis and stochastic attention, the limited independence of the mechanistic validation, and the generator\textquotesingle s extrapolation and privacy behavior. We addressed each concern with new analyses and revised text.

We made four main additions. First, we compared stochastic attention with four alternative generators in the same 18-dimensional principal-component space (R1.1, R1.4). Second, we used a convex-hull analysis to ask whether generated profiles lay inside or beyond the region defined by the real cohort (R2.4). Third, we tested whether a synthetic cohort could reveal that an already-known de-identified profile had been stored in the memory matrix (R1.5). We then traced this membership signal to the sampling distribution and derived that distribution in closed form. Fourth, we tested robustness using populations with known covariance, repeated subsets of the clinical cohort, and a nonlinear simulated population (R1.8, R2.1). These analyses separated recovery of a population from reproduction of one small sample and showed the limits of the linear representation.

Two results changed how we interpreted the method. When the principal-component representation was held constant, the fitted two-component mixture and copula achieved marginal errors and mechanistic overlaps similar to those of stochastic attention. These results were therefore not specific to stochastic attention, although the comparison did not quantify the contribution of principal-component analysis itself. Empirical correlation error and novelty differed across generators. We therefore removed the claim that stochastic attention performed best overall. We instead evaluated covariance recovery in simulations where the population covariance was known.

The privacy analysis found a clear membership signal at the reported inverse-temperature value. Follow-up tests showed that the signal came from the sampling distribution in this small-cohort setting, not from the numerical sampler. Lowering the inverse temperature moved synthetic profiles farther from stored profiles but reduced cross-visit covariance fidelity. We also showed that the stochastic-attention target was exactly a finite Gaussian mixture centered on the stored profiles. It could therefore be sampled directly without Langevin dynamics. Direct samples reproduced the reported results, including the membership signal. The reported synthetic cohort and all clinical results remained unchanged. The test distinguished whether an already-known de-identified profile had been stored in the memory matrix; it did not identify a participant or reconstruct an unknown record.

Several findings reflected the small source cohort and the choices used to represent and condition the data. These included the convex-hull result, the membership signal, the limits of conditional generation, and the modest downstream gain. We therefore framed the clinical study as a proof of concept. We removed claims of equivalence and limited the proposed uses to exploratory simulation and workflow testing. For statistical inference, the sample size remained the number of real patients, not the number of generated profiles.

We also clarified that the mechanistic validation was not independent of the source cohort (R1.2, R2.3). We added limitations concerning irregular visit timing, missing data, tail compression, and the linear representation. We simplified the Hopfield explanation and consolidated its methodological details in the Methods (R1.9, R2.minor2). Finally, we reorganized the figures so that the main text still contained seven figures after the new analyses were added (R2.minor1).

\begin{center}\rule{0.5\linewidth}{0.5pt}\end{center}

\section{Reviewer 1}\label{reviewer-1}

\subsection{Major comments}\label{major-comments}

\subsubsection{R1.1: Disentangle PCA preprocessing from the SA energy landscape (ablation baselines in PCA space)}\label{r11-disentangle-pca-preprocessing-from-the-sa-energy-landscape-ablation-baselines-in-pca-space}

\begin{quote}
The regularization may already arise from the PCA dimensionality-reduction step rather than SA itself. Baselines do not disentangle PCA from the SA energy landscape. Suggests baselines operating in the same PCA-reduced space: (i) PCA+GMM, (ii) PCA+KDE, (iii) PCA+nearest-neighbor interpolation, (iv) PCA+diffusion/interpolation, (v) PCA+copula.
\end{quote}

\emph{We agree that understanding why stochastic attention worked in this setting required separating the PCA representation from the generator.} We therefore held the representation constant and compared stochastic attention with four generators that used the same 18-dimensional principal-component space: a fitted two-component diagonal-covariance Gaussian mixture, a kernel-density estimator, \(k\)-nearest-neighbor interpolation, and a Gaussian copula. We decoded and evaluated all five methods in the same way. The fitted mixture was distinct from the exact stochastic-attention mixture, which had one component centered on each stored profile and was sampled directly only as a separate check.

Holding the representation constant allowed us to determine which outcomes changed with the generator. The nearest-neighbor method produced profiles close to the stored records. Its median novelty was 0.05, where novelty was defined as \(1-\max_k\cos(\hat{\xi},\hat{m}_k)\). Only 5\% of its profiles exceeded the 0.2 cutoff, which required a cosine similarity below 0.8 to every stored profile. Marginal error was the median relative error between generated and real per-feature means. The kernel-density estimator had a marginal error of 1.9\% and a mechanistic overlap of 0.82. An overlap of 0.82 meant that 82\% of its pooled predicted-to-recorded thrombin-generation ratios fell within the 5th to 95th percentile range of the real ratios.

The fitted two-component mixture and copula were competitive with stochastic attention. Their marginal errors were 1.1\% and 1.3\%, compared with 1.2\% for stochastic attention. Their mechanistic overlaps were 0.92 and 0.93, compared with 0.90. All profiles generated by stochastic attention and the copula exceeded the novelty cutoff, as did 94\% of the fitted-mixture profiles. Stochastic attention had a larger cross-visit correlation error than the fitted mixture and copula (0.64, compared with 0.42 for both). However, this metric measured agreement with the correlation matrix of the same 23 patients used to construct the stochastic-attention memory matrix or fit the baseline generators. It assessed in-sample reproduction rather than population recovery. We tested population covariance recovery separately in simulations with a known population covariance (R1.8/R2.1). Stochastic attention had lower error than the sample-covariance Gaussian generator in 35 of 39 low-rank \(n<p\) settings. Holding the PCA representation constant showed that the similar marginal errors and mechanistic overlaps were not specific to stochastic attention, whereas empirical correlation error and novelty differed across generators. The comparison did not quantify the contribution of PCA itself because it did not include a matched non-PCA condition. No method performed best on every measure. In the Discussion, we separated these findings from the sampler-specific choice of inverse temperature.

We did not include a neural diffusion model because estimating one from 23 observations would not provide a credible small-sample baseline. The nearest-neighbor method provided a non-neural interpolation comparison.

\textbf{Changes:} Methods (PCA-space comparison and implementation details); Results (Marginal Plausibility); Discussion (representation and sampler distinction); Supplementary Table S13 (five-generator comparison). Blue, with the shared Discussion transition in violet. \texttt{\textbackslash{}rone\{...\}} \texttt{\textbackslash{}rboth\{...\}}

\begin{center}\rule{0.5\linewidth}{0.5pt}\end{center}

\subsubsection{R1.2: Mechanistic validation is not fully independent (shared cohort/variables/assumptions)}\label{r12-mechanistic-validation-is-not-fully-independent-shared-cohortvariablesassumptions}

\begin{quote}
The BZ2012 model operates on the same coagulation variables used for generation, is calibrated on the same cohort, and shares the same biological domain assumptions. Temper the "mechanistically indistinguishable" claim and clarify the residual coupling.
\end{quote}

\emph{We agree that the original independence claim was too strong.} The mechanistic model was calibrated on the same cohort used to construct the stochastic-attention memory matrix. The analysis therefore tested whether real and synthetic profiles were consistent with the same-cohort model, not whether the synthetic profiles agreed with independent biological evidence. The comparison remained informative because the model was blind to the generation process and applied the same fixed calculation to real and synthetic profiles. We revised the Abstract and Results to make this distinction clear. In the final Discussion paragraph, we now state the dependence, explain that it limits the result to same-model consistency, and identify testing with an independently calibrated model as the way to address it. We did not add that experiment because no suitable independent cohort was available.

\textbf{Changes:} Abstract wording; Results (Mechanistic Consistency) clarification; reorganized final Discussion paragraph. Violet (shared with R2.3). \texttt{\textbackslash{}rboth\{...\}}

\begin{center}\rule{0.5\linewidth}{0.5pt}\end{center}

\subsubsection{R1.3: Statistical reliability of very small subgroups (PCOS n=3, PE n=5)}\label{r13-statistical-reliability-of-very-small-subgroups-pcos-n3-pe-n5}

\begin{quote}
Subgroup distributions may reflect interpolation between a handful of trajectories rather than genuine subgroup variability. Add discussion of reliability, explicit acknowledgment of under-sampling, and caution on downstream clinical interpretation.
\end{quote}

\emph{We agree and made this limitation explicit.} The conditioned cohorts reproduced the patterns observed in the small subgroups, but generating 100 profiles did not create more independent subgroup evidence. Conditioning increased the sampling weights of the three profiles from patients with PCOS or the five profiles from patients who developed preeclampsia. Profile magnitudes were still drawn from all 23 patients. The participation ratios were 4.6 and 7.7. These values described how strongly repeated draws were concentrated around a few stored profiles; they were not independent patient counts. Treating the 100 generated profiles as independent observations would therefore understate uncertainty. We limited the proposed uses to exploratory hypothesis generation, workflow testing, and planning future studies. The Discussion now states that the conditioned cohorts still relied on only those three or five patients and that the generated profiles did not add clinical evidence beyond the source cohort.

\textbf{Changes:} Results (Conditional Generation) caveat; Methods (participation-ratio explanation); Discussion. \texttt{\textbackslash{}rone\{...\}}

\begin{center}\rule{0.5\linewidth}{0.5pt}\end{center}

\subsubsection{R1.4: Stronger and more diverse baselines}\label{r14-stronger-and-more-diverse-baselines}

\begin{quote}
CTGAN performs poorly at small n; TVAE was not configured for concatenated longitudinal generation; MVN is disadvantaged in rank-deficient settings. Recommend stronger baselines designed for low-sample tabular generation or manifold interpolation.
\end{quote}

\emph{We agree that stronger small-sample baselines were needed.} As described in R1.1, we added a fitted two-component diagonal-covariance Gaussian mixture, a kernel-density estimator, nearest-neighbor interpolation, and a Gaussian copula. All four used the same reduced space as stochastic attention. This avoided comparing stochastic attention only with methods that operated in the full-dimensional data or generated each visit separately.

The fitted two-component mixture and copula matched stochastic attention closely on marginal error and mechanistic overlap. Nearest-neighbor interpolation produced mostly near-copies, whereas the kernel-density estimator had larger marginal error and lower mechanistic overlap. All profiles generated by stochastic attention and the copula exceeded the novelty cutoff, as did 94\% of the fitted-mixture profiles. Holding the representation constant showed that the marginal and mechanistic results were not specific to stochastic attention, whereas empirical correlation error and novelty differed across generators. The comparison did not quantify the contribution of principal-component analysis itself, and no generator performed best on every measure.

\textbf{Changes:} Same as R1.1. \texttt{\textbackslash{}rone\{...\}}

\begin{center}\rule{0.5\linewidth}{0.5pt}\end{center}

\subsubsection{R1.5: Privacy and memorization (membership inference, re-identification)}\label{r15-privacy-and-memorization-membership-inference-re-identification}

\begin{quote}
Given K=23 and direct interpolation between stored memory patterns, evaluate membership inference risk and patient re-identification concerns.
\end{quote}

\emph{We agree that privacy and memorization required direct testing.} We asked a specific question: if someone already had a candidate de-identified assay profile, could a synthetic cohort reveal whether that profile had been included in the memory matrix? This was not a test of personal identification. The source data contained no direct identifiers or links to named individuals, and neither the synthetic data nor our analysis identified a participant. When the first analysis showed a clear membership signal, we performed additional tests to determine its source.

To answer this question, we tested whether distance from a candidate profile to the synthetic cohort could distinguish profiles used to construct the memory matrix from profiles withheld from it. We repeated the analysis across 40 random splits. In each split, we repeated the preprocessing with 15 patients, constructed the memory matrix, generated 100 synthetic profiles, and held out the other 8 patients. We assigned each real profile a membership score equal to the negative distance to its nearest synthetic profile, so a higher score provided stronger evidence that the profile had been stored. Stored profiles were closer to the synthetic cohort than held-out profiles (median distances of 10.7 and 15.4). We quantified this discrimination using the area under the receiver operating characteristic curve (AUC). Here, the AUC was the probability that the test ranked a randomly chosen stored profile as more likely to have been stored than a randomly chosen held-out profile. A value of 0.5 meant chance performance, and 1.0 meant perfect separation. The mean AUC was 0.97, and 9 of 40 splits reached 1.0. Thus, the test usually distinguished stored from held-out profiles, but it did not guarantee the correct answer for every record.

We next tested whether the membership signal came from the numerical sampling procedure. We compared the reported Langevin sampler with three alternatives that changed its starting points, discarded early samples and reduced dependence between retained samples, or added a Metropolis correction. The correction acceptance rate was approximately 0.999, which indicated negligible discretization error. A Metropolis-adjusted procedure left the mean AUC nearly unchanged at 0.96, and the other sampler changes had little effect on distance to the closest real record. Lowering the inverse temperature over a 40-fold range increased the median distance from synthetic profiles to their nearest stored profiles, but cross-visit covariance fidelity declined.

We then used the direct sampler for the stochastic-attention target, without a Markov chain. For the unit-normalized profiles used here, the weighted Hopfield energy defined an exact finite Gaussian mixture centered on the stored profiles. The probability of selecting each component was proportional to its multiplicity weight, and inverse temperature controlled the common within-component variance. Direct sampling therefore required only selecting a stored profile according to its weight and adding Gaussian noise. These samples reproduced the reported results: marginal error was 1.7\% compared with 1.3\%, cross-visit covariance error was 0.60 compared with 0.65, mechanistic overlap was 0.91 compared with 0.89, median novelty was 0.50 compared with 0.51, and median distance to the closest record was 13.9 compared with 14.0. Results from 100 independently seeded cohorts showed that this agreement did not depend on one favorable seed. Across the same 40 stored and held-out splits, the mean AUC was 0.98. Because the direct sampler had no initialization, burn-in, or discretization error, this result showed that the signal came from the target distribution rather than the numerical sampler.

Thus, the membership signal arose from the target distribution in this small-cohort setting, where every stored profile was the center of a mixture component. It did not result from the numerical sampler. Under this narrow test, inclusion could be inferred for an already-known de-identified record. The analysis did not attach a name to that record, reconstruct an unknown patient\textquotesingle s measurements, or identify a participant. We reported the full diagnostics in the Results and a new Supplement subsection and limited the Discussion to the distinction between detecting inclusion of an already-known record and identifying or reconstructing a participant.

\textbf{Changes:} Results (one question-driven paragraph ending with the bounded candidate-record conclusion) and Discussion (bounded candidate-record interpretation); new Supplement subsection "Supplementary Sampling and Privacy Results," condensed into two question-driven paragraphs (E3, E3b, four sampling procedures, the inverse-temperature sweep, and direct sampling); Methods (closed-form sampling distribution and the membership score, receiver operating characteristic curve, and AUC calculation); Introduction (finite-mixture explanation and the source of non-identical samples); corresponding revisions to the computational-cost and interpolation language. Blue. \texttt{\textbackslash{}rone\{...\}}

\begin{center}\rule{0.5\linewidth}{0.5pt}\end{center}

\subsubsection{R1.6: Robustness to irregular visit timing, missingness, variable-length trajectories}\label{r16-robustness-to-irregular-visit-timing-missingness-variable-length-trajectories}

\begin{quote}
The concatenation strategy assumes visits are temporally and physiologically aligned across patients. Discuss robustness to irregular timing, incomplete trajectories, and variable-length records.
\end{quote}

\emph{We agree that the current pipeline required every patient to have the same assays at the same aligned times.} Concatenating visits preserved relationships across visits, but the pipeline could not directly use irregular visit times, missing assays or visits, or variable-length records. Imputation, padding, or a sequence-based method could support these data, but each would add assumptions that would need validation in a small cohort. We added this limitation to the Discussion and presented it together with the limits of the linear principal-component representation.

\textbf{Changes:} Discussion (a unified data-representation paragraph covering aligned visits, missingness, and linear structure, followed by an explicit scope conclusion). \texttt{\textbackslash{}rone\{...\}}

\begin{center}\rule{0.5\linewidth}{0.5pt}\end{center}

\subsubsection{R1.7: Clinical importance of tail compression}\label{r17-clinical-importance-of-tail-compression}

\begin{quote}
SA smooths distributional extremes, but rare/extreme states are often the phenomena of interest. Discuss effect on downstream mechanistic modeling, implications for rare-event hypothesis generation, and whether the framework systematically underestimates severe pathology.
\end{quote}

\emph{We agree that preservation of clinically important tails required clearer discussion.} In the descriptive bootstrap Mann-\/-Whitney diagnostic, four of 24 feature-\/-condition pairs fell below the 90\% no-difference-detected criterion: FIX and plasminogen in the PCOS and Developed PE groups. The relative differences between the real and synthetic subgroup means for these pairs were 9-\/-15\%, and variability among the real patients was large. This power-limited diagnostic was not an equivalence test and did not determine whether the lower fractions reflected mean shifts, compressed tails, or limited power. We therefore did not interpret it as evidence that the method systematically compressed tails.

Truncating the principal-component representation and drawing profile magnitudes from only 23 observed values could nevertheless compress the generated tails. If a generated cohort under-represented severe states, it could not establish how often those states occurred or how severe they became. Although synthetic profiles extended beyond the real cohort\textquotesingle s convex hull, that result did not show that they reproduced the true tails of the population. We therefore stated that these cohorts should not be used to estimate extreme risk without additional real data. We also clarified that the main mechanistic-overlap measure used the 5th to 95th percentile range of the real cohort and therefore did not validate tail behavior. The Kolmogorov--Smirnov tests compared the full distributions but were not designed specifically to test the tails.

\textbf{Changes:} Discussion (a question-led tail-preservation paragraph that explains the evidence, clinical concern, and supported conclusion); links to E2. \texttt{\textbackslash{}rone\{...\}}

\begin{center}\rule{0.5\linewidth}{0.5pt}\end{center}

\subsubsection{R1.8 / R2.1: Generalizability beyond a single cohort}\label{r18--r21-generalizability-beyond-a-single-cohort}

\begin{quote}
Validation is on a single K=23 cohort from a single domain. Encourage external validation or explicit limitation; recommend proof-of-concept framing. (R2: a second dataset, or a synthetic benchmark / simulation study demonstrating robustness to varying data characteristics.)
\end{quote}

\emph{We agree that one clinical cohort could not establish generalizability or robustness to changes in sample size, dimension, or nonlinear structure.} We therefore added three complementary robustness tests. The Results first compared stochastic attention with a sample-covariance Gaussian generator on recovery of a known population covariance. They then tested stochastic attention alone with fewer clinical training patients and compared the two generators on nonlinear populations.

The covariance test asked whether stochastic attention or the Gaussian comparator could more accurately recover a known population covariance from a small, high-dimensional sample. We created simulated datasets by drawing profiles from low-rank Gaussian populations with known covariance. We varied the number of profiles in each dataset (\(n=8\), 15, 23, 40, or 80), the number of variables per profile (\(p=30\), 120, or 216), and the latent rank (\(r=3\), 5, or 10). For each setting, we repeated the standardization, PCA, and inverse-temperature selection and constructed a stochastic-attention memory matrix. For comparison, we estimated a Gaussian generator from the dataset\textquotesingle s sample mean and covariance. The Gaussian generator used diagonal jitter for numerical sampling but no covariance shrinkage. Both generators produced 100 profiles. We measured covariance error as the square root of the summed squared entry-by-entry differences between the generated and known covariance matrices, divided by the corresponding value for the known matrix. This relative Frobenius error allowed comparisons across settings and tested recovery of the population covariance rather than reproduction of the sampled profiles. Stochastic attention had lower error in 35 of the 39 rank-deficient \(n<p\) settings. At \(n=23\) and \(p=216\), the Gaussian generator\textquotesingle s error was approximately 1.5 times the stochastic-attention error after averaging across ranks.

The sample-size test evaluated stochastic attention alone as the clinical training set became smaller. We randomly selected 20, 15, 10, or 8 of the 23 patients and repeated the selection ten times at each size. We repeated the stochastic-attention preprocessing and inverse-temperature selection, constructed a new memory matrix from each subset, generated 100 profiles, and compared every generated cohort with the same complete set of 23 real patients. This provided a common reference across subset sizes. As the subset size decreased from 20 to 8, mean cross-visit correlation error increased from 0.63 to 1.11 and the median error in variable means increased from about 2.0\% to 4.9\%. Mechanistic overlap remained near 0.90.

The nonlinearity test asked how increasing curvature affected recovery of a nonlinear pattern. For five replicate samples of 23 points at each curvature and dimension, we constructed a stochastic-attention memory matrix and estimated the same sample-covariance Gaussian comparator; both generators produced 100 profiles. The noisy S-curve represented a one-dimensional simulated path that resembled the letter S because it had two opposing bends. Zero curvature made this path a straight line, and larger values made the bends stronger. We spread the path across 30 or 120 simulated variables while preserving its shape. We measured within-plane distance from the known curve and covariance error against an independent 50,000-profile population reference as curvature increased. Stochastic-attention samples moved farther from the curve as curvature increased. At curvature 0.5, their mean distance was at least twice the straight-line value after averaging across the five replicates and both dimensions; it remained below this threshold at curvature 0.25. Despite this decline, at every nonzero curvature, stochastic-attention samples remained closer to the curve on average than Gaussian samples. However, the Gaussian generator recovered covariance more accurately, so neither generator performed better on both measures.

These tests showed lower covariance error for stochastic attention in most of the tested small, linear, rank-deficient settings. Marginal and correlation fidelity declined with fewer clinical training patients, although mechanistic overlap remained near 0.90. The nonlinear benchmark exposed a trade-off: at nonzero curvature, stochastic-attention samples remained closer to the curve on average, whereas the Gaussian generator recovered covariance more accurately. We also noted that related studies have applied the same geometry-based operating-point rule and multiplicity weighting to protein sequence generation and conditional steering. We cited those studies as prior applications in another domain, not as clinical validation. We separately framed the present clinical study as a proof of concept and stated that we did not test the method in a second clinical cohort.

\textbf{Changes:} New Methods subsection "Robustness Tests," including the construction and numerical details for both generators, and a single-paragraph Results subsection "Robustness to Sample Size, Dimensionality, and Nonlinearity" that explains the three tests in plain language and identifies which generators were evaluated in each one. The Methods define the S-curve in plain language, and its first Results mention links to the full construction. The Results distinguish construction of the stochastic-attention memory matrix from estimation of the Gaussian comparator using the sample mean and covariance, describe the number of profiles in each simulated dataset, state what each test measured, retain the averaging used for the reported covariance-error ratio, and separate the decline in stochastic-attention curve recovery as curvature increased from the comparison with the Gaussian generator. The Results move from population-covariance recovery to the stochastic-attention clinical subset test and then the two-generator nonlinear-structure benchmark; revised main-text figure (\texttt{fig:sim-recovery}) and caption; added cross-domain context in the Introduction and Discussion; explicit proof-of-concept framing in the Abstract, Introduction, and Discussion. The Discussion also stated that the subgroup test did not establish equivalence. Violet. \texttt{\textbackslash{}rboth\{...\}}

\begin{center}\rule{0.5\linewidth}{0.5pt}\end{center}

\subsubsection{R1.9 / R2.minor2: Theory exposition for a biomedical audience}\label{r19--r2minor2-theory-exposition-for-a-biomedical-audience}

\begin{quote}
The entropy-inflection (β) selection, participation-ratio interpretation, and geometric interpretation of the energy landscape are hard for a broad biomedical audience; add intuition or a schematic. (R2: streamline the Hopfield theory to focus on the applied contribution.)
\end{quote}

\emph{We agree that the theory needed a simpler explanation.} We revised the Methods to explain the two roles of β. In the Hopfield update, low β spread attention across many stored profiles and moved toward their weighted average. High β concentrated attention on individual profiles. In direct sampling, β did not change the probability of selecting each unit-normalized profile. Those probabilities were set by the multiplicity weights. Instead, β controlled how widely samples spread around the selected profile.

We chose β* at the point where the attention-entropy curve bent downward most strongly. We calculated this as the most negative finite-difference second derivative of normalized entropy with respect to log β. We defined the participation ratio as the number of equally weighted mixture components that would give the same concentration of sampling probability. It was not an independent patient count. We also introduced each displayed equation with a complete sentence that states what the equation gives and connects it to the preceding mathematical step. We kept the entropy definition, calculation details, explanation, and resulting values together in the Methods. We did not add another figure because the main text already contained seven figures and Reviewer 2 requested figure consolidation.

\textbf{Changes:} Methods (simplified explanation, entropy calculation, connected lead-ins to displayed equations, and a separate concise participation-ratio paragraph). Violet. \texttt{\textbackslash{}rboth\{...\}}

\begin{center}\rule{0.5\linewidth}{0.5pt}\end{center}

\subsubsection{R1.10: Unvalidated feature categories (fibrinolytic, viscoelastic)}\label{r110-unvalidated-feature-categories-fibrinolytic-viscoelastic}

\begin{quote}
Mechanistic validation covered only thrombin generation, not fibrinolytic or viscoelastic features, which are a substantial portion of the feature space. Discuss how unvalidated categories may influence the generated manifold and whether they contribute meaningfully to the energy landscape; expand future validation.
\end{quote}

\emph{We agree and clarified this limitation.} The rotational thromboelastometry and fibrinolytic features were part of the full profile and influenced the generated profiles through their loadings on the retained principal components. Their median relative errors were comparable to those of the coagulation factors. We removed the mechanistic-model aside from the Marginal Plausibility subsection so that the paragraph ends with its main finding. The final Discussion paragraph retains the limitation: the mechanistic model represented thrombin generation, not fibrinolysis or clot viscoelasticity, so the analysis did not establish biological fidelity across the full profile. Testing these features mechanistically would require models that represent those processes.

\textbf{Changes:} Results (Marginal Plausibility) and reorganized final Discussion paragraph. \texttt{\textbackslash{}rone\{...\}}

\begin{center}\rule{0.5\linewidth}{0.5pt}\end{center}

\subsubsection{R1.11: Downstream improvement may reflect smoothing/variance suppression}\label{r111-downstream-improvement-may-reflect-smoothingvariance-suppression}

\begin{quote}
The synthetic-calibrated model\textquotesingle s improvement over the real-calibrated model may partly reflect smoothing/regularization that suppresses physiologically meaningful heterogeneity.
\end{quote}

\emph{We agree and made this interpretation more direct.} The modest downstream gain (0.94× ratio) did not show that synthetic data were superior. The original Results paragraph already stated that the profiles came from the same real cohort and that the comparison was not independent validation. We therefore focused the reviewer-added text on the distinct interpretation: the gain may instead have reflected a more stable calibration objective produced by sampling the represented distribution more densely. We also connected this interpretation to the tail-compression limitation and did not claim that we had established the cause of the numerical gain.

\textbf{Changes:} Results (Downstream Utility) + Discussion. \texttt{\textbackslash{}rone\{...\}}

\begin{center}\rule{0.5\linewidth}{0.5pt}\end{center}

\subsubsection{R1.12: Computational scalability}\label{r112-computational-scalability}

\begin{quote}
Discuss how the framework behaves as cohort size increases, whether the energy landscape becomes unstable at larger K, and how complexity scales with dimensionality.
\end{quote}

\emph{We added a direct description of computational cost.} In the reported Langevin implementation, every sampling step evaluated all stored profiles. The cost of generating one profile therefore scaled as \(O(TKd)\), where \(T\) was the number of steps, \(K\) was the number of stored profiles, and \(d\) was their dimension. The closed-form result removed the sequential factor \(T\). The component probabilities could be calculated once. Each new profile then required selecting one component, drawing a \(d\)-dimensional Gaussian vector, and decoding the profile. We clarified that the inverse transformation and decoding were still required for every generated profile.

The simulation benchmark did not establish how the method would perform in larger clinical cohorts. Its advantage over the sample-covariance Gaussian generator was concentrated in the \(n<p\) settings, and the clinical analyses included only 23 patients. Stability and fidelity at larger cohort sizes remained open questions.

\textbf{Changes:} Discussion, including a brief connection to the simulation benchmark. The scaling paragraph appears immediately before the final limitations paragraph. \texttt{\textbackslash{}rone\{...\}}

\begin{center}\rule{0.5\linewidth}{0.5pt}\end{center}

\subsection{Minor comments}\label{minor-comments}

\subsubsection{R1.M1: Temper strong wording}\label{r1m1-temper-strong-wording}

\begin{quote}
Claims of "clinically useful synthetic cohorts" and "statistically indistinguishable" may be too strong given the sample size.
\end{quote}

\emph{We agree and revised these claims throughout the manuscript.} The subgroup analysis was a descriptive, low-power Mann--Whitney test. A result of \texttt{p\ \textgreater{}\ 0.05} meant only that the test did not detect a difference in that replicate; it did not establish equivalence. We also limited the utility claim to the exploratory modeling uses tested in this proof-of-concept study.

\textbf{Changes:} Abstract, Introduction, and Discussion wording. \texttt{\textbackslash{}rone\{...\}}

\begin{center}\rule{0.5\linewidth}{0.5pt}\end{center}

\section{Reviewer 2}\label{reviewer-2}

\subsection{Major comments}\label{major-comments-1}

\subsubsection{R2.1: Generalizability beyond a single dataset}\label{r21-generalizability-beyond-a-single-dataset}

\begin{quote}
All validation is on a single cohort (K=23). Application to a second dataset, or a synthetic benchmark / simulation study demonstrating robustness to varying data characteristics, would significantly strengthen the manuscript.
\end{quote}

\emph{We agree that validation in one clinical cohort was a major limitation.} As described in \textbf{R1.8}, we added complementary robustness tests. We compared stochastic attention with a sample-covariance Gaussian generator when recovering a known covariance across several sample sizes, dimensions, and latent ranks; evaluated stochastic attention alone as fewer clinical profiles were used to construct its memory matrix; and compared both generators as a simulated population became more curved. We also cited previous protein-sequence studies as examples of applications in another domain, not as clinical validation, and framed the clinical study as a proof of concept. The simulations tested specific data properties, but we did not test the stochastic-attention generator in a second clinical cohort.

\textbf{Changes:} See R1.8. Violet. \texttt{\textbackslash{}rboth\{...\}}

\begin{center}\rule{0.5\linewidth}{0.5pt}\end{center}

\subsubsection{R2.2: Role and limitations of PCA; nonlinear embeddings}\label{r22-role-and-limitations-of-pca-nonlinear-embeddings}

\begin{quote}
Discuss scenarios where PCA may be insufficient; comment on whether nonlinear embeddings could be incorporated and the expected trade-offs between interpretability, fidelity, and flexibility.
\end{quote}

\emph{We agree that the limits of the linear representation needed to be tested and discussed.} In the S-curve simulation, generated profiles moved farther from the known curve as curvature increased. After averaging across the five replicates and both dimensions, the mean stochastic-attention distance reached twice its value for a straight population between curvatures 0.25 and 0.5. At every nonzero curvature, stochastic-attention samples remained closer to the curve on average than Gaussian samples, but the Gaussian generator recovered covariance more accurately. Neither generator performed best on both measures. We limited this conclusion to the tested S-curve and did not generalize it to all nonlinear populations.

We also stated that nonlinear representations could describe curved or branching relationships more accurately. However, they would require additional choices for estimation, interpretation, and decoding that would be difficult to support with 23 patients. The linear representation was a practical choice for this cohort, not evidence that clinical trajectories were generally linear. When this representation was held constant, the fitted two-component mixture and copula achieved marginal errors and mechanistic overlaps similar to those of stochastic attention. These results were not specific to stochastic attention, but the comparison did not quantify the contribution of the representation itself.

\textbf{Changes:} Discussion (limits of the linear principal-component representation), with a connection to E1. Red. \texttt{\textbackslash{}rtwo\{...\}}

\begin{center}\rule{0.5\linewidth}{0.5pt}\end{center}

\subsubsection{R2.3: Mechanistic validation: shared calibration dependence}\label{r23-mechanistic-validation-shared-calibration-dependence}

\begin{quote}
The ODE is calibrated on the same real dataset used to generate the synthetic data, creating partial dependence. Clarify that the validation shows consistency with a model fitted to the same cohort rather than fully independent mechanistic fidelity; if feasible, use independently calibrated parameters or an alternative model.
\end{quote}

\emph{We agree.} As described in \textbf{R1.2}, we stated that the mechanistic model was calibrated on the same cohort used to construct the memory matrix. The analysis therefore tested consistency with a same-cohort model, not agreement with independent biological evidence. We reorganized the final Discussion paragraph so that this dependence, its consequence for interpretation, and the need for an independently calibrated model appear together.

\textbf{Changes:} See R1.2. Violet. \texttt{\textbackslash{}rboth\{...\}}

\begin{center}\rule{0.5\linewidth}{0.5pt}\end{center}

\subsubsection{R2.4: Interpolation vs. extrapolation (convex hull, novel phenotypes)}\label{r24-interpolation-vs-extrapolation-convex-hull-novel-phenotypes}

\begin{quote}
Does SA generate novel out-of-sample variation or primarily interpolate? To what extent do generated patients lie within the convex hull of observed data? Does it produce biologically plausible but previously unseen phenotypes?
\end{quote}

\emph{We addressed this question with geometric, biological, and mechanistic checks.} We first calculated the convex hull, the smallest convex region containing all 23 real profiles in the 18-dimensional principal-component space. All 100 synthetic profiles lay outside this hull. However, each real profile also lay outside the hull formed by the other 22 real profiles. Thus, every real profile was an outer corner of this sparse, high-dimensional region, and an inside-or-outside test alone could not show whether the synthetic profiles extrapolated unusually far.

We therefore compared distances from the hull. Our hypothesis was that excessive extrapolation would place synthetic profiles farther from the real cohort than real profiles treated as unseen. Both groups needed to be evaluated against hulls containing 22 real profiles. The leave-one-out calculation already did this for each real profile. For each synthetic profile, we removed its nearest real profile before constructing the comparison hull. The median distances were 11.0 for synthetic profiles and 11.9 for held-out real profiles. A descriptive Mann--Whitney test did not detect a difference (\(p=0.07\)). This result did not establish equivalence, but it provided no evidence that synthetic profiles lay farther from the observed cohort than held-out real profiles. It also showed why hull membership alone was not informative in this setting. The hull occupied little of the surrounding space, and the generator\textquotesingle s normalization meant that a generated direction could lie inside it only by matching a stored direction exactly.

Distance from the hull did not determine whether a profile was biologically plausible, so we applied two additional checks. The first required values that should be positive to remain positive and required coagulation-factor and thrombin-generation values to remain within five real-cohort standard deviations of the mean. Fifty-four of 100 synthetic profiles passed. The other 46 contained a negative estradiol or progesterone value. The unconstrained multivariate-normal baseline had a similar result, with negative hormone values in 43 of 100 profiles. Neither method enforced positivity. The hormones were closer to zero relative to their observed variation than the coagulation factors, which explained why only hormones crossed zero in these samples. The positive factor values did not provide a general guarantee of positivity.

The second check used the BZ2012 model to test whether the generated coagulation-factor combinations produced thrombin-generation behavior consistent with the real cohort. Estradiol and progesterone were not model inputs, and all nine generated factor inputs were positive. We therefore included all 100 synthetic profiles. For each thrombin-generation feature, 83\% to 92\% of profiles had a model-predicted-to-recorded ratio within the 5th to 95th percentile range of the real cohort. Forty-one profiles met this criterion for every feature and visit, and 24 passed both the biological and strict mechanistic criteria. Positive-variable transformations or constrained decoding could prevent negative values. Together, these checks separated geometric extrapolation from biological plausibility. Hull distance provided no evidence that synthetic profiles extended farther than real profiles treated as unseen, but only 24 of 100 profiles passed both criteria. Thus, hull position alone could not determine whether an individual generated profile was biologically plausible. The Supplement reported the hormone-specific counts, extreme values, and comparison with the real-cohort scale.

\textbf{Changes:} Revised Results paragraph separating geometric extrapolation from biological plausibility and stating the bounded conclusion supported by both checks; Supplement (E2), including a corrected and condensed convex-hull table caption. \texttt{\textbackslash{}rtwo\{...\}}

\begin{center}\rule{0.5\linewidth}{0.5pt}\end{center}

\subsubsection{R2.5: Demonstration of real-world utility}\label{r25-demonstration-of-real-world-utility}

\begin{quote}
Does synthetic data improve predictive models beyond the original cohort? Enhance power for subgroup analyses? Enable concrete new hypotheses? Even a limited quantitative demonstration would strengthen translational significance.
\end{quote}

\emph{We agree that the paper needed a concrete account of its practical value.} The revised Results placed the calibration result alongside the subgroup and factor-level analyses. Together, these results supported exploratory modeling and hypothesis generation. Confirming those hypotheses and establishing clinical value would require additional real patients or independent experimental data.

We also clarified in the Introduction that a further motivation was to support the development of deep learning models. The Discussion now cites related work that predicted TGA responses from coagulation factors and searched for factor changes that produce a target response. Although we did not evaluate either task with profiles generated by stochastic attention, we identify these applications as a clear direction for extending the present work. As noted in R1.11, the 0.94× downstream error ratio did not add independent evidence about the broader population. It may instead have reflected a more stable calibration objective.

\textbf{Changes:} Introduction (original motivating application); Results (Downstream Utility); Discussion (concise forward- and inverse-modeling example). \texttt{\textbackslash{}rtwo\{...\}}

\begin{center}\rule{0.5\linewidth}{0.5pt}\end{center}

\subsection{Minor comments}\label{minor-comments-1}

\subsubsection{R2.minor1: Consolidate figures}\label{r2minor1-consolidate-figures}

\begin{quote}
Several figures could move to the supplementary material to improve readability.
\end{quote}

\emph{We agree.} We moved the principal-component-by-visit comparison of stochastic attention and the multivariate-normal baseline to the Supplement. We placed the simulation benchmark (\texttt{fig:sim-recovery}) in its former position. The main text therefore retained seven figures despite the new analyses. We reported the baseline, convex-hull, membership, and sampling diagnostics in Supplementary tables and figures rather than adding more main-text panels.

\textbf{Changes:} Figure placement (floats.tex / supplementary.tex). \texttt{\textbackslash{}rtwo\{...\}}

\begin{center}\rule{0.5\linewidth}{0.5pt}\end{center}

\subsubsection{R2.minor2: Streamline Hopfield theory exposition}\label{r2minor2-streamline-hopfield-theory-exposition}

\begin{quote}
Some sections (particularly the Hopfield network theory) could be streamlined to focus on the applied contribution.
\end{quote}

As described in \textbf{R1.9}, we consolidated the theory and calculation details in the Methods and added a plain-language explanation.

\textbf{Changes:} See R1.9. Violet. \texttt{\textbackslash{}rboth\{...\}}

\begin{center}\rule{0.5\linewidth}{0.5pt}\end{center}

\subsection{Cross-reference of comments and implemented changes}\label{cross-reference-of-comments-and-implemented-changes}

{\def\LTcaptype{none} % do not increment counter
\begin{longtable}[]{@{}llll@{}}
\toprule\noalign{}
Comment & Type & Vehicle & Color \\
\midrule\noalign{}
\endhead
\bottomrule\noalign{}
\endlastfoot
R1.1, R1.4 & New experiment & E1 PCA-space ablation suite & blue \\
R1.2, R2.3 & Prose (temper) & Mechanistic independence clarification & violet \\
R1.3 & Prose & Subgroup reliability caveats & blue \\
R1.5 & New experiment & E3 membership inference / privacy & blue \\
R1.6 & Prose & Irregular-timing / missingness limitations & blue \\
R1.7 & Prose & Tail-compression clinical importance & blue \\
R1.8, R2.1 & New experiment + framing & Simulation benchmark + cross-domain evidence & violet \\
R1.9, R2.minor2 & Writing & Theory streamlining + corrected intuition & violet \\
R1.10 & Prose & Fibrinolytic/viscoelastic features & blue \\
R1.11 & Prose & Downstream smoothing caveat & blue \\
R1.12 & Prose & Computational scalability & blue \\
R1.M1 & Writing & Temper wording & blue \\
R2.2 & Prose & PCA limits / nonlinear embeddings & red \\
R2.4 & New experiment & E2 convex-hull / extrapolation & red \\
R2.5 & Framing & Downstream utility translational framing & red \\
R2.minor1 & Layout & Figure consolidation & red \\
\end{longtable}
}

\end{document}
